\documentclass{article}

% Match the compact public-bank template used by the existing Logic assessment.
\usepackage[letterpaper,margin=0.3in,includehead,includefoot,headheight=14pt,headsep=6pt]{geometry}
\usepackage[T1]{fontenc}
\usepackage{amsmath,amssymb}
\usepackage{enumitem}
\usepackage{fancyhdr}
\usepackage{microtype}
\usepackage{multicol}
\usepackage[compact,small]{titlesec}
\usepackage[most]{tcolorbox}
\usepackage[hidelinks]{hyperref}

\setlength{\parindent}{0pt}
\setlength{\parskip}{0pt}
\setlength{\columnsep}{0.15in}
\titlespacing*{\section}{0pt}{2pt}{1pt}
\titlespacing*{\subsection}{0pt}{2pt}{1pt}
\setlist{nosep,leftmargin=*}
\setlist[itemize]{noitemsep,topsep=0pt,parsep=0pt,partopsep=0pt,leftmargin=*}
\setlist[enumerate]{noitemsep,topsep=0pt,parsep=0pt,partopsep=0pt,leftmargin=*}
\setlength{\abovedisplayskip}{2pt}
\setlength{\belowdisplayskip}{2pt}
\setlength{\abovedisplayshortskip}{1pt}
\setlength{\belowdisplayshortskip}{1pt}

\pagestyle{fancy}
\fancyhf{}
\fancyhead[L]{\textbf{Logic Practice Bank -- 25-Problem Draft}}
\fancyfoot[C]{\thepage}
\renewcommand{\headrulewidth}{0.4pt}

\newcommand{\problem}[2]{%
  \par
  \subsection*{#1 {\normalfont\footnotesize (#2)}}%
}

\tcbset{
  colback=white,
  colframe=black,
  boxrule=0.4pt,
  arc=2pt,
  left=5pt,right=5pt,top=3pt,bottom=3pt
}

\begin{document}

\textbf{Faculty-review practice bank.} This is not a live Core form. Each
labeled item is one bank problem, including all lettered parts. The five items
in each section have the same structure and comparable workload. Sources refer
to Kenneth Rosen, \textit{Discrete Mathematics and Its Applications}, 7th ed.

\begin{multicols}{3}

\section*{A. Construct and interpret truth tables}

\textbf{Common directions for A.1--A.5.} Construct a complete eight-row truth
table. Show useful intermediate columns and classify the final proposition as
a tautology, contradiction, or contingency.

\problem{A.1}{1.1 Exercise 37(a)}
\[
  p\to(\neg q\lor r)
\]

\problem{A.2}{1.1 Exercise 37(b)}
\[
  \neg p\to(q\to r)
\]

\problem{A.3}{1.1 Exercise 37(c)}
\[
  (p\to q)\lor(\neg p\to r)
\]

\problem{A.4}{1.1 Exercise 37(d)}
\[
  (p\to q)\land(\neg p\to r)
\]

\problem{A.5}{1.1 Exercise 37(e)}
\[
  (p\leftrightarrow q)\lor(\neg q\leftrightarrow r)
\]
\section*{B. Translate meaning and diagnose translation errors}

\textbf{Common directions for B.1--B.5.} Each problem has three parts:
(a) translate conditional language into propositional logic; (b) translate
between English and predicate logic; and (c) diagnose a proposed translation.
Define every predicate and domain not already supplied. For part (c), say
whether the proposal is correct, repair it if needed, and explain the error in
one sentence.

\problem{B.1}{adapted from 1.2 Exercise 1 and 1.5 Exercise 3(b)--(c)}
\begin{enumerate}[label=(\alph*)]
  \item Let $e$ mean ``you can edit a protected encyclopedia entry'' and $a$
        mean ``you are an administrator.'' Translate: ``You cannot edit a
        protected entry unless you are an administrator.''
  \item Let $Q(x,y)$ mean ``student $x$ sent a message to student $y$.''
        Translate $\exists x\,\forall y\,Q(x,y)$ into ordinary English.
  \item A student says the formula in (b) means ``Every student sent a message
        to at least one student.'' Diagnose the claim.
\end{enumerate}

\columnbreak
\problem{B.2}{adapted from 1.2 Exercise 2 and 1.5 Exercise 3(b)--(c)}
\begin{enumerate}[label=(\alph*)]
  \item Let $m$ mean ``you can see the movie,'' $e$ mean ``you are over 18,''
        and $p$ mean ``you have parental permission.'' Translate: ``You can
        see the movie only if you are over 18 or have parental permission.''
  \item With $Q(x,y)$ as in B.1, translate $\forall x\,\exists y\,Q(x,y)$
        into ordinary English.
  \item A student translates ``One student sent a message to every student''
        as $\forall y\,\exists x\,Q(x,y)$. Diagnose the translation.
\end{enumerate}

\problem{B.3}{adapted from 1.2 Exercise 3 and 1.5 Exercise 4(d)}
\begin{enumerate}[label=(\alph*)]
  \item Let $g$ mean ``you can graduate,'' $r$ mean ``you completed your major
        requirements,'' $m$ mean ``you owe the university money,'' and $b$
        mean ``you have an overdue library book.'' Translate: ``You can
        graduate only if you completed the requirements, owe no money, and
        have no overdue library book.''
  \item Let $P(x,y)$ mean ``student $x$ has taken computer-science course
        $y$.'' Translate $\exists y\,\forall x\,P(x,y)$ into English.
  \item A student instead writes $\forall x\,\exists y\,P(x,y)$ for the
        sentence in (b). Diagnose the change in meaning.
\end{enumerate}

\columnbreak
\problem{B.4}{adapted from 1.2 Exercise 4 and 1.5 Exercise 9(b)--(c)}
\begin{enumerate}[label=(\alph*)]
  \item Let $w$ mean ``you may use the airport wireless network,'' $d$ mean
        ``you pay the daily fee,'' and $s$ mean ``you subscribe to the
        service.'' Translate: ``To use the network, you must pay the daily fee
        unless you are a subscriber.''
  \item Let $L(x,y)$ mean ``person $x$ loves person $y$.'' Translate
        ``Everybody loves somebody.''
  \item A student claims that $\exists y\,\forall x\,L(x,y)$ is an equivalent
        translation of (b). Diagnose the claim.
\end{enumerate}

\problem{B.5}{adapted from 1.2 Exercise 7(c) and 1.5 Exercise 11(f)}
\begin{enumerate}[label=(\alph*)]
  \item Let $p$ mean ``the message is scanned for viruses'' and $q$ mean ``the
        message came from an unknown system.'' Translate: ``It is necessary to
        scan a message whenever it comes from an unknown system.''
  \item Let $S(x)$ mean ``$x$ is a student,'' $F(y)$ mean ``$y$ is a faculty
        member,'' and $A(x,y)$ mean ``$x$ asked $y$ a question.'' Translate:
        ``Some student has asked every faculty member a question.''
  \item A student uses $A(y,x)$ in place of $A(x,y)$ in (b). Diagnose what the
        reversed argument order says.
\end{enumerate}
\section*{C. Negate quantified statements}

\textbf{Common directions for C.1--C.5.} For (a), define predicates, symbolize
the statement, negate it with no negation immediately before a quantifier, and
translate the result into ordinary English. For (b), rewrite the given formula
so that every negation appears immediately before a predicate.

\problem{C.1}{adapted from 1.4 Exercise 32(a) and 1.5 Exercise 30(a)}
\begin{enumerate}[label=(\alph*)]
  \item All dogs have fleas. The domain is all things.
  \item $\neg\exists y\,\exists x\,P(x,y)$
\end{enumerate}

\problem{C.2}{adapted from 1.4 Exercise 32(d) and 1.5 Exercise 30(b)}
\begin{enumerate}[label=(\alph*)]
  \item No monkey can speak French. The domain is all things.
  \item $\neg\forall x\,\exists y\,P(x,y)$
\end{enumerate}

\columnbreak
\problem{C.3}{adapted from 1.4 Exercise 33(a) and 1.5 Exercise 33(a)}
\begin{enumerate}[label=(\alph*)]
  \item Some old dogs can learn new tricks. The domain is all things.
  \item $\neg\forall x\,\forall y\,P(x,y)$
\end{enumerate}

\problem{C.4}{adapted from 1.4 Exercise 33(e) and 1.5 Exercise 33(b)}
\begin{enumerate}[label=(\alph*)]
  \item No one in this class knows both French and Russian. The domain is all
        people.
  \item $\neg\forall y\,\exists x\,P(x,y)$
\end{enumerate}

\problem{C.5}{adapted from 1.4 Exercise 34(b) and 1.5 Exercise 33(c)}
\begin{enumerate}[label=(\alph*)]
  \item All Swedish movies are serious. The domain is all things.
  \item $\neg\forall y\,\forall x\bigl(P(x,y)\lor Q(x,y)\bigr)$
\end{enumerate}
\section*{D. Determine consistency and satisfiability}

\textbf{Common directions for D.1--D.5.} Define one atomic proposition for
each simple statement and translate every condition. Determine whether the
conditions are consistent. Give a satisfying assignment when one exists; when
none exists, give a short contradiction. Complete the additional balancing
request in the final sentence of each problem.

\problem{D.1}{adapted from 1.2 Exercise 9}
The system is in multiuser state if and only if it is operating normally. If
the system is operating normally, the kernel is functioning. The kernel is not
functioning or the system is in interrupt mode. If the system is not in
multiuser state, it is in interrupt mode. The system is not in interrupt mode.
If the conditions are inconsistent, identify a minimal collection of the given
conditions that already forces the contradiction.

\problem{D.2}{adapted from 1.2 Exercise 10}
Whenever system software is being upgraded, users cannot access the file
system. If users can access the file system, they can save new files. If users
cannot save new files, the software is not being upgraded. In addition to
deciding consistency, list all satisfying assignments for the three atomic
propositions.

\columnbreak
\problem{D.3}{adapted from 1.2 Exercise 11}
The router can send packets to the edge system only if it supports the new
address space. Supporting the new address space requires the latest software
release. If the latest software release is installed, the router can send
packets to the edge system. The router does not support the new address space.
In addition to deciding consistency, determine which truth values are forced.

\problem{D.4}{adapted from 1.2 Exercise 12}
If the file system is not locked, new messages are queued. The file system is
not locked if and only if the system is functioning normally. If messages are
not queued, they are sent to the message buffer. If the file system is not
locked, messages are sent to the buffer. Messages are not sent to the buffer.
In addition to deciding consistency, explain which conclusions follow from the
last condition.

\problem{D.5}{adapted from 1.2 Exercise 18}
Jasmine becomes unhappy if she attends a party with Samir. Samir attends only
if Kanti attends. Kanti will not attend unless Jasmine attends. Determine which
subsets of Jasmine, Samir, and Kanti can be invited without violating these
conditions. State every satisfying combination, including the possibility of
inviting no one.
\section*{E. Derive logical equivalences with cited rules}

\begin{tcolorbox}[title=Permitted rules for E.1--E.5]
Give a chain of equivalences from the left-hand expression to the right-hand
expression. Cite exactly one rule at every step. Each rule may be used in either
direction.
\begin{itemize}
  \item \textbf{Implication:} $a\to b\equiv\neg a\lor b$.
  \item \textbf{Double Negation:} $\neg\neg a\equiv a$.
  \item \textbf{De Morgan:}
        $\neg(a\lor b)\equiv\neg a\land\neg b$ and
        $\neg(a\land b)\equiv\neg a\lor\neg b$.
  \item \textbf{Distributive:}
        $a\lor(b\land c)\equiv(a\lor b)\land(a\lor c)$ and
        $a\land(b\lor c)\equiv(a\land b)\lor(a\land c)$.
  \item \textbf{Associative:}
        $(a\lor b)\lor c\equiv a\lor(b\lor c)$ and
        $(a\land b)\land c\equiv a\land(b\land c)$.
  \item \textbf{Commutative:}
        $a\lor b\equiv b\lor a$ and $a\land b\equiv b\land a$.
  \item \textbf{Idempotent:}
        $a\lor a\equiv a$ and $a\land a\equiv a$.
\end{itemize}
Changing parentheses requires Associative; changing order requires
Commutative.
\end{tcolorbox}

\problem{E.1}{1.3 Exercise 22}
Show that
\[
  (p\to q)\land(p\to r)
  \quad\text{and}\quad
  p\to(q\land r)
\]
are logically equivalent.

\problem{E.2}{1.3 Exercise 23}
Show that
\[
  (p\to r)\land(q\to r)
  \quad\text{and}\quad
  (p\lor q)\to r
\]
are logically equivalent.

\problem{E.3}{1.3 Exercise 24}
Show that
\[
  (p\to q)\lor(p\to r)
  \quad\text{and}\quad
  p\to(q\lor r)
\]
are logically equivalent.

\problem{E.4}{1.3 Exercise 25}
Show that
\[
  (p\to r)\lor(q\to r)
  \quad\text{and}\quad
  (p\land q)\to r
\]
are logically equivalent.

\problem{E.5}{1.3 Exercise 26}
Show that
\[
  \neg p\to(q\to r)
  \quad\text{and}\quad
  q\to(p\lor r)
\]
are logically equivalent.
\end{multicols}
\end{document}
